\documentclass[12pt,a4paper,openany,oneside]{book}

% --- Paquetes ---
\usepackage[utf8]{inputenc}
\usepackage[spanish, es-noshorthands]{babel}
\usepackage{amsmath, amsfonts, amssymb}
\usepackage{graphicx}
\usepackage{geometry}
\usepackage{hyperref}
\usepackage{booktabs}
\usepackage{fancyhdr}
\usepackage{listings}
\usepackage{xcolor}
\usepackage{tikz}
\usepackage{pgfplots}
\usetikzlibrary{arrows.meta, positioning, shapes.geometric}
\pgfplotsset{compat=1.18}

\geometry{margin=2.5cm}

% --- Colores y estilos para código Python ---
\definecolor{codegreen}{rgb}{0,0.6,0}
\definecolor{codegray}{rgb}{0.5,0.5,0.5}
\definecolor{codepurple}{rgb}{0.58,0,0.82}
\definecolor{backcolour}{rgb}{0.95,0.95,0.92}

\lstset{
    language=Python,
    backgroundcolor=\color{backcolour},   
    commentstyle=\color{codegreen},
    keywordstyle=\color{blue},
    numberstyle=\tiny\color{codegray},
    stringstyle=\color{codepurple},
    basicstyle=\ttfamily\small,
    breaklines=true,
    frame=single,
    showstringspaces=false
}

% --- Estilo de cabecera y pie ---
\pagestyle{fancy}
\fancyhf{}
\renewcommand{\headrulewidth}{0.4pt}
\renewcommand{\footrulewidth}{0.4pt}
\fancyhead[L]{\small Carlos César Sánchez Coronel}
\fancyhead[R]{\small Sesión 6: La Derivada - Midiendo el Cambio}
\fancyfoot[L]{\small Carlos César Sánchez Coronel}
\fancyfoot[C]{\thepage}

% --- Portada ---
\title{
    \textbf{Sesión 6} \\ 
    \Huge \textbf{LA DERIVADA: MIDIENDO EL CAMBIO} \\
    \large De la pendiente instantánea al entrenamiento de redes neuronales \\

}
\author{
    \small Por \\ \vspace{0.2cm}
    \Large \textsc{Carlos César Sánchez Coronel}
}
\date{2026}

\begin{document}

\maketitle
\tableofcontents
\addcontentsline{toc}{chapter}{Índice general}

% ------------------------------------------------------------------
% CAPÍTULO 6: SESIÓN 6 - LA DERIVADA
% ------------------------------------------------------------------
\chapter{La Derivada: Midiendo el Cambio}

En esta sesión abordamos el concepto fundamental de derivada, su interpretación geométrica y física, las reglas de derivación y, lo más importante, su papel central en el entrenamiento de modelos de inteligencia artificial. Comprender la derivada es esencial para entender cómo una máquina ``aprende'' ajustando parámetros mediante el descenso de gradiente y el algoritmo de backpropagation.

\section{Introducción y Motivación}

\subsection{¿Por qué estudiar derivadas?}
La derivada es la herramienta matemática para medir el \textbf{cambio instantáneo}. Sin derivadas, no tendríamos:
\begin{itemize}
    \item \textbf{Ingeniería:} Para calcular la resistencia de materiales ante cargas variables.
    \item \textbf{Economía:} Para hallar el costo marginal y maximizar ganancias.
    \item \textbf{Física:} Para entender la relación entre posición, velocidad y aceleración.
    \item \textbf{Inteligencia Artificial:} Para los algoritmos de optimización que actualizan los pesos de una red neuronal.
\end{itemize}

\subsection{Origen Geométrico: De la Secante a la Tangente}
La derivada nace de la necesidad de medir la inclinación de una curva en un punto exacto. Históricamente, esto se resuelve mediante el paso de una recta secante a una recta tangente.

\begin{itemize}
    \item \textbf{Recta Secante:} Línea que corta a la función en dos puntos, $(x_1, y_1)$ y $(x_2, y_2)$. Su pendiente es:
    \[
    m_{\text{sec}} = \frac{\Delta y}{\Delta x} = \frac{f(x_2) - f(x_1)}{x_2 - x_1}
    \]
    \item \textbf{Recta Tangente:} A medida que el segundo punto se acerca al primero ($\Delta x \to 0$), la secante se convierte en tangente. La pendiente de esta tangente es la \textbf{derivada}.
\end{itemize}

\begin{center}
\begin{tikzpicture}
\begin{axis}[
    title={Aproximación de la tangente mediante secantes},
    xlabel=$x$,
    ylabel=$f(x)$,
    xmin=0, xmax=3,
    ymin=0, ymax=4,
    grid=major,
    samples=100,
    domain=0:3
]
\addplot[blue, thick] {0.5*x^2 + 1};
\addplot[red, thick, domain=0.5:2.5] {1.5*x - 0.5};
\addplot[green, thick, domain=0.8:2.2] {1.2*x + 0.2};
\addplot[orange, thick, domain=1:2] {1*x + 1};
\node[pin=90:{$f(x)$}] at (axis cs:2.5,3.5) {};
\node[pin=90:{Secante}] at (axis cs:1.8,2.5) {};
\node[pin=90:{Secante más cercana}] at (axis cs:1.5,2.2) {};
\node[pin=90:{Tangente}] at (axis cs:1.2,2.2) {};
\end{axis}
\end{tikzpicture}
\captionof{figure}{Aproximación de la recta tangente mediante secantes cada vez más cercanas al punto de interés.}
\end{center}

\subsection{El Método Delta ($\Delta$)}
El método delta es la definición formal de la derivada mediante límites. Representamos un incremento pequeño en $x$ como $\Delta x$:
\begin{equation}
\boxed{ \frac{dy}{dx} = \lim_{\Delta x \to 0} \frac{f(x + \Delta x) - f(x)}{\Delta x} }
\end{equation}
Esta expresión nos da la \textbf{pendiente instantánea} de la función en cualquier punto $x$.

\section{Notación y Aplicaciones}

\subsection{Sistemas de Notación y Orden Superior}
La diferenciación puede aplicarse más de una vez sobre una función. Dependiendo del autor y del campo de estudio (Ingeniería, Física o Matemática Pura), se utilizan distintas formas de representar las derivadas de orden superior.

\begin{table}[h]
\centering
\begin{tabular}{@{}lllll@{}}
\toprule
\textbf{Orden} & \textbf{Leibniz} & \textbf{Lagrange} & \textbf{Euler} & \textbf{Newton} \\ \midrule
1ª Derivada & $\frac{dy}{dx}$ & $f'(x)$ & $D_x y$ & $\dot{y}$ \\ \addlinespace
2ª Derivada & $\frac{d^2y}{dx^2}$ & $f''(x)$ & $D_x^2 y$ & $\ddot{y}$ \\ \addlinespace
3ª Derivada & $\frac{d^3y}{dx^3}$ & $f'''(x)$ & $D_x^3 y$ & $\dddot{y}$ \\ \addlinespace
$n$-ésima Derivada & $\frac{d^ny}{dx^n}$ & $f^{(n)}(x)$ & $D_x^n y$ & - \\ \bottomrule
\end{tabular}
\caption{Comparativa de notaciones para derivadas de orden superior.}
\end{table}

\subsection{La Derivada en Ecuaciones (¿Cómo leerlas?)}
Para notar la diferencia visual y práctica, observa cómo se presenta la misma relación física (la aceleración de un objeto) en distintas notaciones dentro de una ecuación:

\begin{enumerate}
    \item \textbf{Forma de Lagrange (Matemática):} $f''(x) = a$
    \item \textbf{Forma de Leibniz (Ingeniería):} $\frac{d^2s}{dt^2} = a$
    \item \textbf{Forma de Newton (Física):} $\ddot{s} = a$
\end{enumerate}

\subsection{Aplicación Geométrica: La Recta Tangente}
La derivada $f'(a)$ nos da la pendiente $m$ de la recta tangente a la curva en el punto $a$. Con esto, podemos construir la ecuación de la recta:
\[
y - f(a) = f'(a)(x - a)
\]
Que simplificada toma la forma clásica $y = mx + b$.

\subsection{Ejemplo de Orden Superior}
Si tenemos $f(x) = x^4$, sus derivadas sucesivas son:
\begin{align*}
f'(x) &= 4x^3 \quad &\text{(Velocidad de cambio)}\\
f''(x) &= 12x^2 \quad &\text{(Aceleración de cambio)}\\
f'''(x) &= 24x \\
f^{(4)}(x) &= 24 \quad &\text{(Constante)}
\end{align*}

\section{Tabla de Derivadas y Reglas}

\subsection{Tabla de Derivadas más Usadas}
A continuación, se presentan las derivadas de las funciones más comunes:

\begin{table}[h]
\centering
\begin{tabular}{ll}
\toprule
\textbf{Función $f(x)$} & \textbf{Derivada $f'(x)$} \\ \midrule
Constante $k$ & $0$ \\
Identidad $x$ & $1$ \\
Potencia $x^n$ & $n \cdot x^{n-1}$ \\
Exponencial $e^x$ & $e^x$ \\
Exponencial $a^x$ & $a^x \ln a$ \\
Logaritmo natural $\ln x$ & $\frac{1}{x}$ \\
Logaritmo $\log_a x$ & $\frac{1}{x \ln a}$ \\
Seno $\sin x$ & $\cos x$ \\
Coseno $\cos x$ & $-\sin x$ \\
Tangente $\tan x$ & $\sec^2 x$ \\
Arcoseno $\arcsin x$ & $\frac{1}{\sqrt{1-x^2}}$ \\
Arcocoseno $\arccos x$ & $-\frac{1}{\sqrt{1-x^2}}$ \\
Arcotangente $\arctan x$ & $\frac{1}{1+x^2}$ \\ \bottomrule
\end{tabular}
\caption{Tabla resumen de derivadas elementales.}
\end{table}

\subsection{Reglas de Derivación}
\begin{itemize}
    \item \textbf{Múltiplo constante:} $[c \cdot f(x)]' = c \cdot f'(x)$.
    \item \textbf{Suma o Resta:} $(f \pm g)' = f' \pm g'$.
    \item \textbf{Producto:} $(f \cdot g)' = f'g + fg'$.
    \item \textbf{Cociente:} $\left( \frac{f}{g} \right)' = \frac{f'g - fg'}{g^2}$.
    \item \textbf{Regla de la Cadena:} $\frac{d}{dx}[f(g(x))] = f'(g(x)) \cdot g'(x)$.
\end{itemize}

\section{Ejercicios Resueltos Paso a Paso}

\subsection{Caso Polinómico y Raíces}
Derivar $f(x) = 4x^5 - 2x^2 + \sqrt{x}$. 

\textit{Solución:} Re-escribimos la raíz como potencia: $f(x) = 4x^5 - 2x^2 + x^{1/2}$.
\begin{align*}
f'(x) &= (4 \cdot 5)x^4 - (2 \cdot 2)x^1 + \frac{1}{2}x^{-1/2} \\
&= 20x^4 - 4x + \frac{1}{2\sqrt{x}}
\end{align*}

\subsection{Caso con Regla del Producto}
Derivar $f(x) = x^2 \cdot \cos x$. 

\textit{Solución:} 
\begin{align*}
f'(x) &= (x^2)'\cos x + x^2(\cos x)' \\
&= 2x\cos x - x^2\sin x
\end{align*}

\subsection{Caso con Regla del Cociente}
Derivar $f(x) = \frac{x^2 + 1}{x^3 - x}$.

\textit{Solución:} Aplicamos la regla del cociente:
\begin{align*}
f'(x) &= \frac{(2x)(x^3 - x) - (x^2 + 1)(3x^2 - 1)}{(x^3 - x)^2} \\
&= \frac{2x^4 - 2x^2 - (3x^4 + 3x^2 - x^2 - 1)}{(x^3 - x)^2} \\
&= \frac{2x^4 - 2x^2 - 3x^4 - 2x^2 + 1}{(x^3 - x)^2} \\
&= \frac{-x^4 - 4x^2 + 1}{(x^3 - x)^2}
\end{align*}

\subsection{Caso con Regla de la Cadena}
Derivar $f(x) = \ln(x^3 + 1)$. 

\textit{Solución:} Aplicamos la derivada del logaritmo y luego la cadena:
\[
f'(x) = \frac{1}{x^3 + 1} \cdot (x^3 + 1)' = \frac{3x^2}{x^3 + 1}
\]

\subsection{Caso de la Cadena con Potencia}
Derivar $f(x) = (3x^2 + 1)^4$:
\[
f'(x) = 4(3x^2 + 1)^3 \cdot (6x) = 24x(3x^2 + 1)^3
\]

\section{Optimización: Máximos y Mínimos}

Antes de ver cómo aprende una IA, debemos entender cómo las derivadas nos permiten encontrar los mejores valores posibles para una función. A esto le llamamos \textbf{Optimización}.

\subsection{Criterio de la Primera Derivada}
Para encontrar un punto donde la función alcanza un pico (máximo) o un valle (mínimo), buscamos los \textbf{Puntos Críticos}. Estos ocurren cuando la pendiente es horizontal, es decir:
\[
f'(x) = 0
\]

\begin{itemize}
    \item \textbf{Máximo Relativo:} La función sube y luego baja. La pendiente pasa de positiva a negativa.
    \item \textbf{Mínimo Relativo:} La función baja y luego sube. La pendiente pasa de negativa a positiva.
\end{itemize}

\subsection{Criterio de la Segunda Derivada}
¿Cómo sabemos si el punto donde $f'(x)=0$ es un máximo o un mínimo sin ver la gráfica? Usamos la segunda derivada:
\begin{enumerate}
    \item Si $f''(x) > 0$, la función es "cóncava hacia arriba" (curvatura positiva) $\rightarrow$ Es un \textbf{Mínimo}.
    \item Si $f''(x) < 0$, la función es "cóncava hacia abajo" (curvatura negativa) $\rightarrow$ Es un \textbf{Máximo}.
\end{enumerate}

\subsection{Guía Paso a Paso para el Análisis}
Para resolver cualquier problema de optimización en la evaluación, sigue este protocolo:
\begin{enumerate}
    \item \textbf{Derivar:} Obtén la primera derivada $f'(x)$.
    \item \textbf{Puntos Críticos:} Iguala $f'(x) = 0$ y resuelve para $x$.
    \item \textbf{Segunda Derivada:} Calcula $f''(x)$.
    \item \textbf{Clasificar:} Sustituye tu punto crítico en $f''(x)$ y observa el signo.
\end{enumerate}

\subsection{Ejercicio Resuelto: Maximizando Beneficios}
\textit{Problema:} Una empresa de software estima que sus ganancias están dadas por $G(x) = -x^2 + 8x - 12$. Hallar el precio óptimo $x$ para maximizar la ganancia.

\textbf{Paso 1: Hallar la primera derivada.}
\[
G'(x) = -2x + 8
\]

\textbf{Paso 2: Encontrar el punto crítico.}
Igualamos a cero:
\[
-2x + 8 = 0 \implies 2x = 8 \implies x = 4
\]

\textbf{Paso 3: Verificación con la segunda derivada.}
\[
G''(x) = -2
\]
Como $G''(4) = -2 < 0$, la función es cóncava hacia abajo. Por lo tanto, en \textbf{$x = 4$} existe un \textbf{Máximo Relativo}.

\subsection{Conexión Crítica con la IA: ¿Por qué optimizamos?}
Imagina que estás enseñando a una IA a predecir el precio de una casa. Si la casa vale \$100k y la IA dice \$120k, hay un error de \$20k. La herramienta que mide ese error es la \textbf{Función de Pérdida} (Loss Function).

\subsubsection{La ecuación del error: Error Cuadrático Medio (MSE)}
Piensa en la función de pérdida más básica:
\[
L(\hat{y}, y) = (\hat{y} - y)^2
\]
Donde:
\begin{itemize}
    \item $L$: es el valor de la pérdida (Loss).
    \item $\hat{y}$: es lo que la IA \textbf{predice}.
    \item $y$: es el valor \textbf{real} que debió obtener.
\end{itemize}
El exponente al cuadrado se usa para que el error siempre sea positivo y para penalizar más fuerte los errores grandes.

\subsubsection{Interpretación para el entrenamiento}
\begin{itemize}
    \item \textbf{Pérdida Alta:} La predicción está muy lejos de la realidad. El modelo está "confundido".
    \item \textbf{Pérdida Baja (Mínimo):} La predicción es casi idéntica al valor real. El modelo ha "aprendido".
\end{itemize}

\subsubsection{Derivada de la función MSE}
Si tenemos un modelo lineal $\hat{y} = wx + b$, la pérdida para un conjunto de datos es:
\[
L(w,b) = \frac{1}{N} \sum_{i=1}^{N} (wx_i + b - y_i)^2
\]
Las derivadas parciales respecto a $w$ y $b$ nos indican cómo ajustar estos parámetros para minimizar el error. Por ejemplo:
\[
\frac{\partial L}{\partial w} = \frac{2}{N} \sum_{i=1}^{N} (wx_i + b - y_i) x_i
\]
Esta derivada es la base del \textbf{descenso de gradiente}.

\section{Derivadas en IA: Funciones de Activación}

El aprendizaje profundo (\textit{Deep Learning}) se fundamenta en la optimización de funciones de coste mediante el descenso del gradiente. Para ello, es imprescindible calcular derivadas de las funciones que componen la red, especialmente las \textbf{funciones de activación}. Estas introducen no linealidad, permitiendo que la red aprenda patrones complejos.

\subsection{Función Sigmoide: La puerta a la probabilidad}

\begin{equation}
\sigma(x) = \frac{1}{1 + e^{-x}}
\end{equation}

\begin{itemize}
    \item \textbf{Dominio:} $\mathbb{R}$, \textbf{Rango:} $(0,1)$.
    \item \textbf{Continuidad y derivabilidad:} Infinitamente diferenciable.
\end{itemize}

\subsubsection{Derivada de la sigmoide}
\begin{equation}
\sigma'(x) = \sigma(x) \bigl(1 - \sigma(x)\bigr)
\end{equation}

\textbf{Demostración paso a paso:}
\begin{align*}
\sigma(x) &= (1 + e^{-x})^{-1} \\
\sigma'(x) &= -1 (1 + e^{-x})^{-2} \cdot (-e^{-x}) = \frac{e^{-x}}{(1 + e^{-x})^2} \\
&= \frac{1}{1 + e^{-x}} \cdot \frac{e^{-x}}{1 + e^{-x}} = \sigma(x) \cdot \frac{e^{-x}}{1 + e^{-x}}.
\end{align*}
Pero $\frac{e^{-x}}{1 + e^{-x}} = 1 - \frac{1}{1 + e^{-x}} = 1 - \sigma(x)$. Luego $\sigma'(x) = \sigma(x)(1 - \sigma(x))$.

\subsubsection{Casos de uso en IA}
\begin{itemize}
    \item Clasificación binaria (capa final).
    \item Puertas en redes LSTM.
\end{itemize}

\subsection{Función ReLU: El motor del Deep Learning moderno}

\begin{equation}
\text{ReLU}(x) = \max(0, x)
\end{equation}

\begin{itemize}
    \item \textbf{Dominio:} $\mathbb{R}$, \textbf{Rango:} $[0, \infty)$.
    \item \textbf{Continuidad:} Continua; no derivable en $x=0$.
\end{itemize}

\subsubsection{Derivada de ReLU}
\[
\text{ReLU}'(x) = 
\begin{cases}
1, & x > 0 \\
0, & x \le 0
\end{cases}
\]
En la práctica, se define como 0 en $x=0$ (o a veces 0.5).

\subsubsection{Casos de uso en IA}
\begin{itemize}
    \item Capas ocultas de redes convolucionales (CNN).
    \item Redes profundas en general.
\end{itemize}

\subsection{Función Softplus: Versión suave y diferenciable}

\begin{equation}
\text{Softplus}(x) = \ln(1 + e^x)
\end{equation}

\begin{itemize}
    \item \textbf{Dominio:} $\mathbb{R}$, \textbf{Rango:} $(0, \infty)$.
    \item \textbf{Derivabilidad:} Infinitamente diferenciable.
\end{itemize}

\subsubsection{Derivada de Softplus}
\[
\text{Softplus}'(x) = \sigma(x) = \frac{1}{1 + e^{-x}}
\]

\subsubsection{Casos de uso}
\begin{itemize}
    \item Alternativa suave a ReLU cuando se requiere derivada definida en todo punto.
\end{itemize}

\subsection{Función GELU: El secreto de los Transformers}

\begin{equation}
\text{GELU}(x) = x \cdot \Phi(x)
\end{equation}
donde $\Phi(x)$ es la función de distribución acumulada normal estándar. En la práctica se usa la aproximación:
\[
\text{GELU}(x) \approx 0.5x \left(1 + \tanh\left[\sqrt{\frac{2}{\pi}}\left(x + 0.044715 x^3\right)\right]\right)
\]

\begin{itemize}
    \item \textbf{Dominio y rango:} $\mathbb{R} \to \mathbb{R}$.
    \item \textbf{Derivabilidad:} Suave.
\end{itemize}

\subsubsection{Casos de uso}
\begin{itemize}
    \item Modelos Transformer (BERT, GPT).
\end{itemize}

\section{Python: Cálculo de Derivadas con SymPy y Derivación Numérica}

\subsection{Cálculo simbólico con SymPy}
SymPy permite obtener derivadas exactas de funciones simbólicas.

\begin{lstlisting}[language=Python]
import sympy as sp
import numpy as np
import matplotlib.pyplot as plt

# Definir variable simbólica
x = sp.Symbol('x')

# Derivada de una función polinómica
f = x**3 + 2*x**2 - 5*x + 1
f_prima = sp.diff(f, x)
print("f'(x) =", f_prima)

# Derivada de una función compuesta
g = sp.sin(x**2)
g_prima = sp.diff(g, x)
print("g'(x) =", g_prima)

# Evaluar la derivada en un punto
valor = f_prima.subs(x, 2)
print("f'(2) =", valor)

# Derivada de la sigmoide
sig = 1 / (1 + sp.exp(-x))
sig_prima = sp.diff(sig, x)
print("sig'(x) =", sig_prima)
\end{lstlisting}

\subsection{Derivación numérica y comparación con la analítica}
Podemos aproximar derivadas numéricamente usando diferencias finitas y comparar con el valor exacto.

\begin{lstlisting}[language=Python]
def f(x):
    return x**3

def f_prima_exacta(x):
    return 3*x**2

def derivada_numerica(f, x, h=1e-5):
    return (f(x + h) - f(x)) / h

x_val = 2.0
h_vals = [1e-1, 1e-2, 1e-3, 1e-4, 1e-5, 1e-6]
print("Exacta:", f_prima_exacta(x_val))
for h in h_vals:
    aprox = derivada_numerica(f, x_val, h)
    error = abs(aprox - f_prima_exacta(x_val))
    print(f"h={h:.0e}: {aprox:.10f}, error={error:.2e}")
\end{lstlisting}

\subsection{Visualización de la derivada}
\begin{lstlisting}[language=Python]
x_vals = np.linspace(-2, 2, 100)
y_vals = f(x_vals)
y_prima_vals = f_prima_exacta(x_vals)

plt.figure(figsize=(10,4))
plt.subplot(1,2,1)
plt.plot(x_vals, y_vals, label='f(x)')
plt.title('Función')
plt.grid(True)

plt.subplot(1,2,2)
plt.plot(x_vals, y_prima_vals, label="f'(x)", color='red')
plt.title('Derivada')
plt.grid(True)
plt.show()
\end{lstlisting}

\section{Notación en Papers Científicos}
En artículos de IA, las derivadas aparecen en contextos como:
\begin{itemize}
    \item Definición de gradiente: $\nabla_{\mathbf{w}} \mathcal{L}(\mathbf{w}) = \left( \frac{\partial \mathcal{L}}{\partial w_1}, \dots, \frac{\partial \mathcal{L}}{\partial w_n} \right)$.
    \item Actualización de pesos: $\mathbf{w}_{t+1} = \mathbf{w}_t - \eta \nabla \mathcal{L}(\mathbf{w}_t)$.
    \item Regla de la cadena en backpropagation: $\frac{\partial \mathcal{L}}{\partial w_{ij}^{(l)}} = \frac{\partial \mathcal{L}}{\partial a_j^{(l)}} \cdot \frac{\partial a_j^{(l)}}{\partial z_j^{(l)}} \cdot \frac{\partial z_j^{(l)}}{\partial w_{ij}^{(l)}}$.
\end{itemize}

\section{Ejercicios Propuestos}
\begin{enumerate}
    \item Calcula la derivada de $f(x) = \ln(\sin x)$ usando la regla de la cadena.
    \item Obtén la derivada de la función de pérdida de entropía cruzada binaria $L = -y \log \hat{y} - (1-y) \log(1-\hat{y})$ respecto a $\hat{y}$.
    \item Implementa en Python la derivada numérica de $f(x) = e^{-x^2}$ y compárala con la analítica.
    \item Demuestra que la derivada de la función $\tanh$ es $1 - \tanh^2(x)$.
    \item Encuentra los puntos críticos de $f(x) = x^3 - 3x$ y clasifícalos usando la segunda derivada.
\end{enumerate}

\section{Resumen}
En esta sesión hemos aprendido:
\begin{itemize}
    \item La definición de derivada como límite de la pendiente secante.
    \item Las reglas de derivación (potencia, producto, cociente, cadena).
    \item La interpretación geométrica y su aplicación en optimización.
    \item El papel crucial de las derivadas en el entrenamiento de modelos de IA (funciones de pérdida, funciones de activación).
    \item Cómo calcular derivadas simbólica y numéricamente con Python.
\end{itemize}
La derivada es la base del descenso de gradiente y del backpropagation, temas que abordaremos en sesiones posteriores.

\end{document}